\documentclass[10pt,a4paper]{article}
\usepackage{amsmath,amssymb,amsthm,graphicx}
\usepackage{amsmath,amssymb,amsfonts,amsthm,latexsym,graphicx,enumitem}
\usepackage{hyperref}
\usepackage{thmtools,thm-restate}
\usepackage[capitalise,nameinlink]{cleveref}
\setcounter{MaxMatrixCols}{15}
\newcommand{\N}{\mathbb{N}}
\newcommand{\Q}{\mathbb{Q}}
\newcommand{\R}{\mathbb{R}}
\newcommand{\Z}{\mathbb{Z}}
\newcommand{\F}{\mathbb{F}}
\newcommand{\zero}{\mathbf{0}}
\newcommand{\Id}{\mathbf{I}}
\newcommand{\vecu}{\vec{u}}
\newcommand{\Acal}{\mathcal{A}}
\newcommand{\Ccal}{\mathcal{C}}
\newcommand{\Hcal}{\mathcal{H}}
\newcommand{\Lcal}{\mathcal{L}}
\newcommand{\Pcal}{\mathcal{P}}
\newcommand{\Rcal}{\mathcal{R}}
\newcommand{\Mcal}{\mathcal{M}}
\newcommand{\M}{\mathcal{M}}
\newcommand{\Scal}{\mathcal{S}}
\newcommand{\Tcal}{\mathcal{T}}
\newcommand{\Fcal}{\mathcal{F}}
\newcommand{\Gcal}{\mathcal{G}}
\newcommand{\Tr}{\mathrm{Tr}}
\newcommand{\OR}{\mathrm{OR}}
\newcommand{\MAJ}{\mathrm{MAJ}}
\newcommand{\symvec}{\mathrm{symvec}}
\newcommand{\symmat}{\mathrm{Sym}}
\newcommand{\MINCUT}{\mathrm{MINCUT}}
\newcommand{\ARGMINCUT}{\mathrm{ARGMINCUT}}
\newcommand{\lin}{\mathrm{lin}}
\newcommand{\cut}{\mathrm{cut}}
\newcommand{\ind}{\mathrm{ind}}
\newcommand{\diag}{\mathrm{diag}}
\newcommand{\IP}{\mathrm{IP}}
\newcommand{\PARITY}{\mathrm{PARITY}}
\newcommand{\upto}{\mathbin{:}}
\DeclareMathOperator*{\argmin}{argmin}
\newcommand{\polylog}{\mathrm{polylog}}
\newcommand{\zeros}{\mathrm{zeros}}
\newcommand{\ones}{\mathrm{ones}}
\newcommand{\rmvec}{\mathrm{vec}}
\newcommand{\rmmat}{\mathrm{mat}}
\newcommand{\laspan}{\mathrm{span}}
\newcommand{\cross}{\mathrm{overlap}}
\newcommand{\coisa}{weighted biadjacency matrix }
\newcommand{\card}[1]{|#1|}
\newcommand{\len}[1]{l(#1)}
\newcommand{\floor}[1]{\lfloor #1 \rfloor}
\newcommand{\ceil}[1]{\left\lceil #1 \right\rceil}
\newcommand{\pair}[2]{\langle #1,#2 \rangle}
\newcommand{\triple}[3]{\langle #1,#2,#3 \rangle}
\def\01{\{0,1\}}

\newcommand{\ith}{i^{\scriptsize \mbox{{\rm th}}}}
\newcommand{\jth}{j^{\scriptsize \mbox{{\rm th}}}}

\newcommand{\cd}{\mathrm{cdim}}
\newcommand{\sep}{\mathrm{sep}}
\newcommand{\mer}{\mathrm{mer}}

\DeclareMathOperator*{\argmax}{arg\,max}

\newcommand{\norm}[1]{\|#1\|}


\newcommand{\troy}[1]{\textcolor{red}{#1}}
\newcommand{\miklos}[1]{\textcolor{blue}{#1}}
\newcommand{\shengyu}[1]{\textcolor{orange}{#1}}
\newcommand{\tongyang}[1]{\textcolor{violet}{#1}}

\newcommand{\bra}[1]{\langle#1|}
\newcommand{\ket}[1]{|#1\rangle}
\newcommand{\bigket}[1]{\big|#1\big\rangle}
\newcommand{\braket}[2]{\langle#1, #2\rangle}
\newcommand{\rk}{\mathrm{rk}}
\newcommand{\I}{\mathcal{I}}
\newcommand{\1}{\mathbf{1}}


\newtheorem{theorem}{Theorem}
\newtheorem{question}[theorem]{Question}
\newtheorem{lemma}[theorem]{Lemma}
\newtheorem{corollary}[theorem]{Corollary}
\newtheorem{proposition}[theorem]{Proposition}
\newtheorem{remark}[theorem]{Remark}
\newtheorem{example}[theorem]{Example}
\newtheorem{claim}[theorem]{Claim}
\newtheorem{fact}[theorem]{Fact}
\newtheorem{hypothesis}[theorem]{Hypothesis}

\theoremstyle{definition}
\newtheorem{exercise}[theorem]{Exercise}
\newtheorem{problem}[theorem]{Problem}
\newtheorem{definition}[theorem]{Definition}
\title{The sharp maximum cut dimension of an n-vertex weighted undirected graph is 2n minus 3}
\author{Troy Lee, Tongyang Li, Miklos Santha, Shengyu Zhang}
\date{doi:10.4230/LIPIcs.CCC.2021.15; CC BY-NC-SA 4.0}
\begin{document}
\maketitle
\noindent\textit{Editorial note.} Complete native CC BY-NC-SA4.0 author version used under its actual license for noncommercial research delivery; NC/SA terms retained. Source comments and unselected alternate proof/material blanked with original locator lines retained. MainTheorems1/2 counted as one sharp result, including complete Jain appendix and explicit attainment construction. Original chapter and global theorem offsets restored. Standard article fonts replace unused Times/class cosmetics; original mathematics and graph model unchanged. Direct linear-query consequence retained only as context. All production/validation solely on Yupeng.\par\medskip
\noindent\textit{Allowed standard prerequisites (editorial).}\begin{itemize}
\item Weighted undirected graphs, cuts, linear span, laminar families and rooted tree representations. Complete original graph/edge-slot/weight/crossing definitions and source-local uncrossing/tree/linear-independence proofs retained; Jain lemma included with full author appendix proof
\end{itemize}\medskip
\section{Introduction}
Let $G = (V,w)$ be a weighted undirected $n$-vertex graph where $w$ is an $\binom{n}{2}$-dimensional nonnegative real vector assigning a (possibly zero) weight to each edge slot.  For a
nontrivial subset $\emptyset \ne X \subsetneq V$, let $\Delta(X)$ be the set of edges of $G$ with one endpoint in $X$ and one endpoint in $\bar{X} = V \setminus X$.  A \emph{cut} $S$ in $G$ is a subset of edges of the form $\Delta(X)$ for
a nontrivial set $X$.  The sets $X$ and $\bar{X}$ are called the \emph{shores} of the cut.  For a cut $S$, its \emph{weight} is the sum of the weights of the edges in $S$, denoted $w(S)$.
The minimum cut problem is to find the minimum of $w(S)$ over all cuts $S$.  The study of algorithms for the minimum cut
problem in theoretical computer science goes back at least to the 1960's and has given rise to a vast and beautiful literature.
Minimum cut is also a problem of great practical importance with applications
to, for example, clustering algorithms and evaluating network reliability. Randomized algorithms can solve the minimum cut problem in nearly linear time: in 1996 Karger gave an
algorithm with running time $O(m \log^3(n))$ to compute the minimum cut of a weighted graph with $m$ edges \cite{Karger00}.  This was the best known bound until very recently when two independent
works improved on it.  Gawrychowski, Mozes, and Weimann \cite{GMW20} gave a randomized algorithm with running time $O(m\log^2(n))$ \cite{GMW20} and Mukhopadhyay and Nanongkai \cite{MN20} gave a randomized algorithm with time complexity $O(m\frac{\log^2(n)}{\log\log n}+n\log^6(n))$.   Gawrychowski, Mozes, and Weimann \cite{GMW20b} later improved the running time of the Mukhopadhyay and Nanongkai algorithm to $O(m\frac{\log^2(n)}{\log\log n}+n\log^{3+\varepsilon}(n))$.

For \emph{simple} graphs $G$, randomized algorithms are known with running times $O(m \log(n))$ and $O(m + n \log^3(n))$ \cite{GNT20}.  For simple graphs even nearly linear
time \emph{deterministic} algorithms are known.  Kawarabayashi and Thorup gave an $O(m \log^{12}(n))$ time algorithm \cite{KT19}, which was subsequently improved to
$O(m (\log(n) \log \log n)^2)$ by Henzinger, Rao, and Wang \cite{HRW20}.


Our work spans two aspects of the study of the minimum cut problem.  The first is to query complexity lower bounds on minimum cut.
A natural model in which to study the query
complexity of minimum cut is for algorithms allowed to make \emph{cut queries}.  A cut query algorithm can query any subset $\emptyset \ne X \subsetneq V$ and receives the
answer $w(\Delta(X))$.  One motivation to study cut query algorithms comes from submodular function minimization.  The cut function $f(X) = w(\Delta(X))$ is a submodular
function, and finding the minimum cut value is equivalent to finding the minimum value of $f$ over all nontrivial sets $X$.  The problem of minimizing a submodular function is often studied with respect to an evaluation
oracle, which in the case of the cut function is exactly a cut query.

Harvey \cite{Harvey08} observed that results on the deterministic communication complexity of deciding graph connectivity \cite{HMT88} imply that any deterministic cut query algorithm to compute minimum cut,
or even to decide if the graph is connected or not, must make at least $c n$ cut queries, for a constant $c < 1$.  Analogous results on the randomized communication complexity of connectivity \cite{BFS86}
imply an $\Omega(n/\log(n))$ lower bound on the number of cut queries needed by a randomized algorithm to compute minimum cut (or even connectivity).

On the algorithms side, Rubinstein, Shramm, and Weinberg \cite{RSW18} gave a randomized algorithm computing the minimum cut of a simple graph with $O(n \log(n)^3)$ many
cut queries\footnote{\cite{RSW18} state the bound as $\tilde O(n)$ but we estimate their sparsifier based algorithm to make $O(n \log(n)^3)$ many cut queries.}.
Recently, Mukhopadhyay and Nanongkai \cite{MN20} used a different approach based on Karger's 2-respecting tree algorithm
\cite{Karger00} to also give a randomized $\tilde O(n)$ cut query algorithm to compute minimum cut in a general undirected weighted graph.

For deterministic cut query algorithms, there remains a large gap between the best upper and lower bounds.  We are not aware of any deterministic algorithm for minimum cut
better than learning the entire graph, which can take $\Omega(n^2/\log(n))$ cut queries in the worst case.  On the lower bound side, Graur, Pollner, Ramaswamy, and Weinberg \cite{GPRW20} recently introduced
a very interesting lower bound technique called the \emph{cut dimension}, which we now describe.  Let $G = (V,w)$ be a weighted undirected graph with $n$ vertices and $m$ edges, and let $\M(G)$ be
the set of minimum cuts of $G$.  For a cut $S \in \M(G)$,
let $\chi(S) \in \{0,1\}^{m}$ be the characteristic vector of $S$
amongst the $m$ edges of $G$.  Let $\vec{\M}(G) = \{\chi(S): S \in \M(G) \}$.  The cut dimension of $G$, denoted
$\cd(G)$, is the dimension of $\laspan(\vec{\M}(G))$.
It is shown in \cite{GPRW20} that for any $n$-vertex graph $G$, the cut dimension $\cd(G)$ is a lower bound on the deterministic cut query complexity of computing
minimum cut on weighted $n$-vertex graphs.  Moreover, for every $n \ge 2$ they construct an $n$-vertex graph $G$ with cut dimension $3n/2-2$.

Besides showing lower bounds on cut query complexity, the cut dimension is a natural measure of the complexity of mincuts in a graph.  There is a rich
literature on the possible structure of mincuts in a graph.  Perhaps the first result of this kind is the cactus representation of mincuts by \cite{DKL76}.
A cactus for a graph $G$ is a sparse weighted graph $C$ that represents all the mincuts of $G$.  One consequence of the
cactus representation is that the number of possible mincuts in an $n$-vertex weighted graph is at most $\binom{n}{2}$.  This upper bound was later given
an algorithmic proof via Karger's famous contraction algorithm (\cite{Karger93}, Theorem~6.1).  The
$n$-vertex cycle graph has $\binom{n}{2}$ many minimum cuts and shows that this bound can be tight.

While the cycle has $\binom{n}{2}$ many mincuts, these cuts live in an $n$-dimensional space as the $n$-vertex cycle only has $n$ edges.  Is it possible to construct graphs
with many cuts that also have high cut dimension?  We show that this is not possible, and in fact the cut dimension of an $n$-vertex graph is at most $2n-3$.

\begin{theorem}[Main Upper Bound]
\label{thm:cd_upper}
For any weighted undirected graph $G$ on $n\ge 2$ vertices it holds that $\cd(G) \le 2n-3$.
\end{theorem}
Like the cactus representation, this shows another aspect in which the mincuts of a graph are constrained to have a relatively simple structure.  We further show that this bound is
tight by constructing graphs with cut dimension $2n-3$ for every $n \ge 2$.
\begin{theorem}[Main Lower Bound]
\label{thm:cd_lower}
For every $n \ge 2$ there exists an $n$-vertex weighted undirected graph $G$ with $\cd(G) = 2n-3$.
\end{theorem}




In addition to shedding further light on the structure of minimum cuts, this improves the best known lower bound on the deterministic cut query complexity of
the minimum cut problem to $2n-3$.  We additionally show that the cut dimension is even a lower bound on a stronger query model called the linear query model, recently
studied in \cite{ACK20}.  In the linear query model, the algorithm can query any vector $x \in \R^{\binom{n}{2}}$ and receives the answer $\braket{w}{x}$, the
inner product of $w$ and $x$.  Linear queries can be much more powerful than cut queries as one can completely learn an unweighted graph with a single linear query.
By an information theoretic argument learning an unweighted graph can require $\Omega(n^2/\log(n))$ many cut queries since each cut query reveals at most
$O(\log (n))$ bits.


\setcounter{theorem}{3}\section{Preliminaries}\label{sec:prelim}

For every natural number $n$, we denote by $[n]$ the set
$\{1,2, \ldots, n\}$.
For a vector $z \in \R^n$ we write $z \geq \mathbf{0}$ if every coordinate of the vector
is at least 0, and similarly we write $z = \mathbf{0}$ if $z$ is the all-zero vector.
We denote the scalar product of two vectors $z,z' \in \R^n$
by $\braket{{z}}{z'}$. For any matrix, denote the rank of $A$ by $\rk(A)$.
We denote the disjoint union of sets $X$ and $Y$ by $X \sqcup Y.$

\subsection{Graphs, cuts, sets}

An undirected {\em weighted graph} on $n$ vertices is a couple
$G = (V,w)$, where $V$ is the set of vertices with $|V| = n$,
the set of edge slots
$V^{(2)}$ is the
set of subsets of $V$ with cardinality 2, and the weight function
$w : V^{(2)} \rightarrow \R $ is non-negative. We refer to the vertex set of $G$ as $V(G)$.
The set of edges of $G$
is defined as $E = \{e \in V^{(2)} : w(e) > 0\}$.
When in a graph $G=(V,w)$ the weight of every edge is 1, we say that the graph is {\em unweighted},
and we refer to it also as $G=(V,E)$; such graph is also called a {\em simple graph}.
For an edge $e = \{u,v\} $, we say that $u$ and $v$ are the endpoints of $e$.
For a subset $X \subseteq V$ of the vertices, we denote by $E(X)$ the set of edges in $E$
which have both endpoints in $X$, and for disjoint subsets $X, Y \subseteq V$, we denote by
$E(X, Y)$ the set of edges with exactly one endpoint in each of the two sets.
We extend the weight function $w$ to any subset $E'$ of the edges by $w(E') = \sum_{e \in E'} w(e)$.
We will deal only with graphs which have at least $2$ vertices.

We fix an ordering
$v_1 < v_2 < \cdots < v_n$ of the vertices which induces also an ordering $\{v_1, v_2\}, \{v_1, v_3\}, \ldots, \{v_{n-1},v_n\}$ of the edge slots as well as an ordering $e_1 < e_2 < \ldots < e_m$ of the $m=|E|$ edges.  We view $w \in \R^{\binom{n}{2}}$ as a vector whose $\ith$ coordinate gives the (possibly zero) weight of the $\ith$ edge slot
according to this ordering,
and we define $\vec{w} \in \{0,1\}^m$ as the restriction of $w$ to the edges.
With some slight abuse of notation, for a set of edges $S \subseteq E$, we use the same symbol $S$ to also denote
the characteristic vector in $\{0,1\}^{\binom{n}{2}}$ of $S$ among all edge slots.
We further need the characteristic vector of $S \subseteq E$ among the $m$ edges $E$, for which we use the
notation $\chi(S) \in \{0,1\}^m$.
For a family ${\cal F}$ of subsets of the edges, we use the notation
$\vec{{\cal F}} = \{ \chi(S) \in \{0,1\}^m : S \in {\cal F}\}$.


For $X \subseteq V$, we denote by $\bar{X}$ the set $V \setminus X$.
A {\em cut} $S$ is a set $E(X, \bar{X})$ for some $\emptyset \neq X \subsetneq V$.
We call
$X$ and $\bar{X}$ the {\em shores} of $S$, and we denote the cut by $\Delta(X)$.
A cut is a {\em star cut} if one of its shores is a singleton,
otherwise it is {\em non-star} cut. If the singleton shore of a star cut $S$ is $\{v\}$, then we say that $S$ is a star cut
at $v$.  The {\em weight} of a cut is the sum of the weights of its edges.
For a cut $S$ we define the {\em graph of the cut} $S$ as the unweighted graph $G(S)= (V', E')$
where $V'$ is the set of vertices in $V$
that are endpoints of at least one edge in $S$, and $E'=S$.
We say that a cut $S$ is
{\em connected} if $G(S)$ is a connected graph.
A cut is a {\em minimum} cut, or mincut, for short, if no other cut has smaller weight.
We denote by $\M(G)$ be the set of minimum cuts of $G$.
The {\em cut dimension} of $G$ is $\cd(G) = \dim(\laspan(\vec{\cal{M}}(G)))$.

Let $V$ be a set of size $n$. Two sets $X,Y \subseteq V$ are said to \emph{overlap} if $X \cap Y \ne \emptyset, \bar{X} \cap Y \ne \emptyset,
X \cap \bar{Y} \ne \emptyset$.
 A family $\Gcal$ of subsets of $V$ is said to be \emph{laminar} if for all $X,Y \in \Gcal$
it holds that $X$ and $Y$ do not overlap. A set family $\Gcal \subseteq 2^{V}$ is said to be \emph{closed under overlaps}
if for every $X,Y \in \Gcal$ that overlap it holds
that $X \cap Y, X \cup Y \in \Gcal$.  A laminar subset $\Lcal \subseteq \Gcal$ is said to be \emph{maximal in} $\Gcal$ if for every $X \in \Gcal - \Lcal$
there is a $Y \in \Lcal$ such that $X,Y$ overlap.  We say a laminar subset $\Lcal$ is maximal if it is maximal in $2^{V}$.


The sets $X,Y \subseteq V$ \emph{cross} if they overlap and additionally $\bar{X} \cap \bar{Y} \ne \emptyset$.
Note that if $X,Y$ cross then so do $X,\bar{Y}$. A set family $\Gcal \subseteq 2^{V}$ is said to be {\em cross-free} if for all $X,Y \in \Gcal$ it holds
that $X$ and $Y$ do not cross. Observe that if $X$ and $Y$ do not cross
then either $Y$ or $\bar{Y}$ is a subset of $X$ or $\bar{X}$.
Let $G = (V,w)$ be a graph with $n$ vertices.
Two cuts $\Delta(X)$ and $\Delta(Y)$ of $G$ are {\em crossing} if $X$ and $Y$ are crossing.
Let $\Fcal = \{ \Delta(X_1), \ldots,  \Delta(X_k) \}$ be a set of cuts of $G$.
We say that $\Fcal$ is {\em cross-free family of cuts} if
$\Gcal = \{X_1, \ldots, X_k\}$ is cross-free.  Note that it does not matter which shore we take to be in $\Gcal$.

There is a close relationship between cross-free families of cuts and laminar sets.
Let $\Fcal = \{ \Delta(X_1), \ldots,  \Delta(X_k) \}$ be a cross-free family of cuts where each $X_i \subseteq V$,
and let $X_i' = X_i$ if $v_1 \not \in X_i$ and $X_i' = \bar{X}_i$ otherwise.
The \emph{beach} of $\Fcal$ is the set $\Gcal = \{X_1', \ldots, X_k'\}$.
For a family of sets $\Gcal \subseteq 2^{V}$ we say that it is \emph{proper} if $\emptyset, V \not \in \Gcal$,
and we say that it is \emph{complement free} if it does not contain $X,Y$ with $Y = \bar{X}$.

\begin{claim}\label{claim:cross-free}
Let $\Fcal$ be a cross-free family of distinct cuts and $\Gcal$ its beach.  Then $\Gcal$ is proper, complement free and laminar.
\end{claim}

\begin{proof}
First, $\Gcal$ does not contain $\emptyset$ or $V$ because these are not shores of cuts.  It is complement free because $\Fcal$ contains
distinct cuts, and its beach contains exactly one representative shore from each cut.  Finally, we show that it is laminar.  Let
$X_1, X_2 \in \Gcal$.  By definition of a beach, neither of these sets contain $v_1$, thus $\bar{X}_1 \cap \bar{X}_2 \ne \emptyset$.  Therefore if
$X_1, X_2$ overlapped they would also cross, in contradiction to $\Fcal$ being a cross-free family of cuts.
\end{proof}

A mincut is {\em crossless} if no other
mincut crosses it. Observe that a star mincut is always crossless.
Also, if a  mincut $\Delta(X)$ is crossless then
for every mincut $\Delta(Y)$, either $Y$ or $\bar{Y}$ is a subset of $X$ or $\bar{X}$.
Crossing mincuts have a nice structural property which was already observed by~\cite{DKL76}.

\begin{claim}
\label{claim:mincuts}
Let $G=(V,w)$ be a weighted graph.
If $\Delta(X), \Delta(Y) \in \M(G)$ cross then
$\Delta(X \cap Y), \Delta(X \cup Y)  \in \M(G)$.
\end{claim}

\begin{proof}
We have $\Delta(X \cap Y) \neq \emptyset $ and $ \Delta(X \cup Y) \neq V$ because $\Delta(X)$ and $\Delta(Y)$ cross. The cut function is submodular therefore we have
\begin{equation*}
\label{eq:mincuts}
w(\Delta(X \cap Y) ) + w(\Delta(X \cup Y) ) \leq w(\Delta(X)) + w(\Delta(Y)).
\end{equation*}
Let $c$ be the weight of a minimum cut in $G$. Then the right hand side of the above inequality
is equal to $2c$, while its left hand side is at least $2c$. Therefore
$w(\Delta(X \cap Y) ) + w(\Delta(X \cup Y) ) = 2c$ from which the statement follows.
\end{proof}


\setcounter{section}{3}\setcounter{theorem}{15}\section{The cut dimension is at most $2n-3$}\label{sec:upper}
In this section we prove \cref{thm:cd_upper} that $\cd(G) \le 2n-3$ for any undirected weighted graph $G$ on $n \ge 2$ vertices.  This will follow from
two facts:
\begin{enumerate}
  \item For $n \ge 2$ a cross-free family of cuts in an $n$-vertex graph has cardinality at most $2n-3$.
  \item If $\Lcal \subseteq \M(G)$ is a maximal cross-free subset of the mincuts of $G$
  then $\laspan(\vec{\Lcal}) = \laspan(\vec{\M}(G))$.
\end{enumerate}
We remind the reader that $\vec{\Lcal} = \{\chi(S) : S \in \Lcal\}$ where $\chi(S) \in \{0,1\}^{|E|}$ is the characteristic
vector of the cut $S$ amongst the edges of $G$.

These two facts are presented in the next two subsections.

\subsection{Cardinality of a cross-free family of cuts}
\label{sec:card}
Recall from \cref{claim:cross-free} that if $\Lcal$ is a cross-free family of cuts then the beach $\Gcal$ of $\Lcal$ is a laminar family of sets.
A standard inductive proof shows that a laminar family of subsets of a universe of cardinality $n$ that contains no singletons has size at most $n-1$, and thus a laminar
family in general has size at most $2n-1$.  A beach has the additional properties of being proper and complement free
which allows one to prove an upper bound of $2n-3$.  This is mentioned by \cite{Goemans06} in the paragraph after Theorem 4
under the heading ``Size of a Laminar Family'', who observes that the standard inductive proof also
implies the bound is attained only if the family includes the universe and at least one set and its complement.
See also Corollary~2.15 of \cite{KV18}, where it is shown that a proper laminar family has cardinality at most $2n-2$.

\begin{lemma}
\label{lem:laminar_size}
Let $n \ge 2$, $V$ a set of cardinality $n$, and $\Gcal \subseteq 2^{V}$ be a family of sets which is proper and laminar.  Then
$|\Gcal| \le 2n-2$.  If $\Gcal$ is proper, laminar, and complement free then $|\Gcal| \le 2n-3$.
\end{lemma}

\begin{proof}
First we show the $2n-2$ upper bound.
We prove by induction.  Consider first the base case where $n=2$ and $V = \{v_1, v_2\}$.  As $\emptyset, V \not \in \Gcal$ the only possible elements to include in $\Gcal$
are $\{v_1\}, \{v_2\}$ and $|\Gcal| \le 2 = 2n-2$.

Now we assume the statement is true for families of sets on a universe of $n-1$ elements and show it holds for families of sets on a universe of size $n$.  Let
$\Gcal \subseteq 2^{V}$ be a proper laminar family.  We say that $X \in \Gcal$ is maximal
if there is no set $Y \in \Gcal$ with $X \subset Y$.  Let $X_1, \ldots, X_m$ be the maximal sets in $\Gcal$.  Note that we must have
$X_i \cap X_j = \emptyset$ for all $i\ne j \in [m]$.  This is because for distinct maximal sets $X_i - X_j, X_j - X_i \ne \emptyset$ thus if
$X_i \cap X_j \ne \emptyset$ they would be overlapping.  If $\cup_{i=1}^m X_i \subsetneq V$ then the result already holds by the induction hypothesis.
Thus we may assume $m \ge 2$ and $X_1, \ldots, X_m$ form a partition of $V$.
The family $\Fcal_1 = \{Y : Y \subsetneq X_1\}$ is a laminar family on the universe $X_1$ which does not contain $X_1$.  Hence by the induction hypothesis it has at most $2|X_1|-2$ many sets.  This holds for all
$i =1, \ldots, m$, thus including $X_1, \ldots, X_m$ the total number of sets is $\sum_{i=1}^m 2|X_i| - m \le 2n-2$.

Now we show the $2n-3$ upper bound additionally assuming the family is complement free.  We show this result directly using the upper bound
of $2n-2$ we have just shown on the size of proper laminar families.  Let $\Gcal \subseteq 2^{V}$ be proper, laminar, and complement free,
and let $X_1, \ldots, X_m$ be the maximal sets in $\Gcal$, which again must be disjoint.  The number of subsets strictly contained in $X_i$ is at most $2|X_i| - 2$ by the previous result.  Thus, including $X_1, \ldots, X_m$
we can upper bound the size of $\Gcal$ by $\sum_{i=1}^m 2|X_i| - m$.  If $m > 2$ then the upper bound of $2n-3$ already holds.  If $m=1$ then as $\Gcal$ is
a proper family we must have $|X_1| \le n-1$ in which case the upper bound of $2n-3$ holds as well.  Finally, consider the case $m=2$.  In this case, if $|X_1 \cup X_2| < n$ then
the bound already holds.  If $X_1 \cup X_2 = V$ then $X_2 = \bar{X}_1$ and we must exclude one of these sets, giving a bound of $2n-2-1 = 2n-3$.
\end{proof}

\begin{remark}
From the proof in the proper, laminar, complement-free case we can observe for what maximal sets equality in the upper bound can hold.  The first is the case where there are three maximal sets
$X_1, X_2, X_3$ that form a partition of $[n]$.  With $V=[6]$ an example of this type saturating the bound is
$\Gcal = \{ \{1\}, \ldots, \{6\}, \{1,2\}, \{3,4\}, \{5,6\} \}$.  The second is the case where there are two maximal sets $X_1,X_2$ that form a
partition of $[n]$ and exactly one of $X_1,X_2$ is not included.  The latter includes the case where there is a single maximal set $X_1$ of size $|X_1| = n-1$.
For $V=[6]$, an example of this type is $\Gcal = \{\{2\}, \ldots, \{6\}, \{2,3\},\{2,3,4\},\{2,3,4,5\},\{2,3,4,5,6\}\}$.
\end{remark}

Chandran and Ram (Lemma~2.13 in \cite{CR04}) show that if the set $\Mcal(G)$ of minimum cuts of a
graph $G$ is cross-free, then $|\Mcal (G)| \le 2n-3$.  This is an
easy corollary of \cref{lem:laminar_size}, which gives something more general.

\begin{corollary}
\label{cor:cross-free_size}
Let $G = (V,w)$ be a graph on $n \ge 2$ vertices.
Let $\Lcal \subseteq \Mcal(G)$ be a subset of minimum cuts that is cross-free.  Then
$|\Lcal| \le 2n-3$.
\end{corollary}

\subsection{Spanning}
Let $\Lcal \subseteq \Mcal(G)$ be a maximal cross-free subset of $\M(G)$.  Here maximal means that for any cut $S \in \M(G) \setminus \Lcal$ there
is a cut $T \in \Lcal$ that crosses $S$.  The fact that $\laspan(\vec{\Lcal}) = \laspan(\vec{\Mcal}(G))$
essentially follows from a key lemma of Jain in his factor of 2 approximation algorithm for the survivable network design problem (Lemma~4.2 in \cite{Jain01}).
Another application of a similar lemma can be found in Goeman's approximation algorithm for the bounded-degree minimum spanning tree problem \cite{Goemans06}.

The context of Jain's lemma is slightly different than ours, as we now explain.  Instead of mincuts, Jain considers the set of cuts
$\Tcal$ which saturate the inequalities of a particular linear program.  He shows that the set $\Tcal$ has the property that if $\Delta(X),\Delta(Y) \in \Tcal$
cross then either
\begin{enumerate}
\item $\Delta(X \cap Y), \Delta(X \cup Y) \in \Tcal$ and $\chi(\Delta(X)) + \chi(\Delta(Y)) = \chi(\Delta(X \cap Y)) + \chi(\Delta(X \cup Y))$, or
\item $X \setminus Y, Y \setminus X \in \Tcal$ and $\chi(\Delta(X)) + \chi(\Delta(Y)) = \chi(\Delta(X \setminus Y)) + \chi(\Delta(Y \setminus X)).$
\end{enumerate}
As shown by Dinitz, Karzanov, and Lomonosov \cite{DKL76}, for crossing mincuts $\Delta(X),\Delta(Y)$ \emph{both} items (1), (2) hold (see \cref{prop:bixby} for a proof).
Thus Jain's lemma applies to $\M(G)$ as well.
\begin{restatable}[\cite{Jain01}]{lemma}{jainspan}\label{thm:span}
Let $G = (V,w)$ be a graph and $\Lcal \subseteq \M(G)$ be a maximal cross-free family of mincuts.
Then $\laspan(\vec{\Lcal}) = \laspan(\vec{\M}(G))$.
\end{restatable}
For completeness, we include a full proof of \cref{thm:span} in \cref{sec:appendix}.

We now can give the first proof of our main upper bound that for any $n\ge 2$ an $n$-vertex graph $G=(V,w)$ has $\cd(G) \le 2n-3$.

\begin{proof}[Proof of \cref{thm:cd_upper}]
Follows from \cref{cor:cross-free_size} and \cref{thm:span}.
\end{proof}

%

\section{Explicit construction of graphs with cut dimension $2n-3$}\label{sec:construction}
In this section we prove \cref{thm:cd_lower} by giving a general technique to explicitly construct graphs of cut dimension $2n-3$.  We focus
on constructing graphs $G = (V,w)$ where $w$ is strictly positive, i.e.\ where $G$ is a complete weighted graph.
The main lemma of this section, \cref{prop:lin_ind}, shows that, in a complete weighted graph, for any cross-free family of cuts $\Lcal$
the vectors in $\vec{\Lcal}$ are linearly independent.

Thus to construct a graph with cut dimension $2n-3$ it suffices to construct a complete weighted graph whose set of mincuts
is a cross-free family of cuts of cardinality $2n-3$.  Such a graph is constructed for every $n \ge 2$ in Theorem~5.2 of \cite{CR04}.
Combining this construction with our linear independence result \cref{prop:lin_ind} gives a proof of our main lower bound \cref{thm:cd_lower}.


In \cref{subsec:construction} we go further and show for any maximal cross-free family $\Fcal \subseteq 2^{[n]}$ there is a complete weighted graph $G = ([n],w)$
with $\M(G) = \{ \Delta(X) : X \in \Fcal\}$.  Moreover, we give an explicit formula for the weight vector $w$.
Part of this construction is a lemma, \cref{lem:min_value},
which may be of independent interest: it says that if $\Lcal$ is a maximal family of cross-free cuts in a graph $G$, and all cuts in $\Lcal$ have the same weight $c$, then $c$
is the weight of the minimum cut in $G$.

A key tool for showing the linear independence of cuts from a cross-free family is the tree representation of a laminar family, which we go over next.

\subsection{Tree representation}
\begin{definition}
For an unweighted directed graph $G=(V,E)$ we let $\delta^+(X) = \{(x,y) \in E : x \in X, y \in V - X\}$.  For a singleton $v \in V$
we write $\delta^+(v)$ instead of $\delta^+(\{v\})$.
\end{definition}

\begin{definition}[Arborescence]
An \emph{arborescence} is a directed rooted tree where all edges point away from the root.  A vertex of an arborescence which is
not the root or a leaf we call an \emph{internal vertex}.
\end{definition}

\begin{definition}[Tree representation]\label{def:tree-repre}
Let $T$ be a directed graph whose underlying undirected graph is a tree.  Let $U$ be a finite set and $\phi: U \rightarrow V(T)$.
For $e = (x,y) \in E(T)$ define $S_e$ as
\[
S_e = \{ s \in U : \phi(s) \mbox{ is in the same connected component of } T - e \mbox{ as } y \} \enspace .
\]
Then $(T,\phi)$ defines a set family $\Fcal = \Fcal(T, \phi)$ where $\Fcal = \{S_e : e \in E(T)\}$.
We say that $(T,\phi)$ is a \emph{tree representation} of $(U, \Fcal)$.
We call $(T,\phi)$ a \emph{faithful tree representation} if $|E(T)| = |\Fcal|$.
For $v \in V(T)$, if there is a $u \in U$ such that $\phi(u) = v$ then we say that $v$ has a \emph{label}.
\end{definition}

We will need the fact that a laminar set family has a faithful tree representation by an arborescence.  A textbook proof of this fact can be found
in Korte and Vygen Proposition 2.14 \cite{KV18}.  While they do not explicitly say the tree representation they construct is faithful,
this is clear from the proof.
\begin{proposition}
\label{prop:treerep}
Let $(U,\Fcal)$ be laminar family.  Then there is a faithful tree representation $(T,\phi)$ of $(U,\Fcal)$ where $T$ is an arborescence.
\end{proposition}

Recall from \cref{claim:cross-free} that if $\Lcal$ is a cross-free family of cuts then its beach $\Gcal$ is laminar, and thus has a tree representation.
\begin{lemma}[Tree structure of maximal cross-free families]
\label{lem:tree_structure}
Let $\Lcal$ be a maximal family of cross-free cuts of a graph $G= ([n],w)$ and $\Gcal \subseteq 2^{[n]}$ its beach.  Then in a faithful tree representation $(T,\phi)$ of $\Gcal$
it holds that
\begin{enumerate}
\item The root $r$ is labeled by $1$ and has $|\delta^+(r)| = 1$
\item There are $n-1$ leaves of $T$ each with a distinct label in $\{2, \ldots, n\}$.
\item Every internal vertex $v$ has $|\delta^+(v)| = 2$.
\end{enumerate}
\end{lemma}

\begin{proof}
As by the definition of a beach, sets do not contain $1$, this means that $1$ must be the label of the root.
As star cuts do not cross any other cut, if $\Lcal$ is maximal it must contain all the star cuts.  This means
that $\Gcal$ contains the sets $\{2\}, \ldots, \{n\}, \{2,\ldots, n\}$.  Thus the outdegree of the root must be $1$, as
this outgoing edge represents the set $\{2,\ldots, n\}$.  Further there must be $n-1$ leaves which are labeled by $2, \ldots, n$.
We have now accounted for all the labels, thus no internal vertex has a label. Further, if there was a leaf $v$ with parent $u$
such that $v$ did not have a label, then $(u,v)$ would represent the
empty set, which by definition is not in $\Gcal$.  Thus there are exactly $n-1$ leaves.

It remains to show that every internal vertex $v$ of $T$ which is not the root has $|\delta^+(v)| = 2$.  Let $v$ be an internal vertex, and as $v$ is not
the root, let $u$ be its parent, and as $v$ is not a leaf let $w$ be a child of $v$.  If $|\delta^+(v)| = 1$ then the edges $(u,v),(v,w)$ would represent the same set,
as $v$ is not labeled.  This contradicts the fact that $(T,\phi)$ is a faithful tree representation.  Now suppose $|\delta^+(v)| > 2$ and let $w,x,y$ be three of its children.  Consider
the sets $X_1, X_2,X_3 \in \Gcal$ represented by the edges $(v,w),(v,x),(v,y)$.  Further the edge $(u,v)$ represents a set $A \in \Gcal$ with
$X_1 \cup X_2 \cup X_3 \subseteq A$.
We claim that in this case $\Lcal$ is not maximal because the cut $\Delta(X_1 \cup X_2)$ does not cross any cut in $\Lcal$.  Indeed, $X_1 \cup X_2$ is contained
in all the sets represented by edges on the path from $v$ to the root, and is disjoint from the sets represented by any other edge of $T$.  Thus we have a contradiction.
\end{proof}

\begin{corollary}
\label{cor:maxcard}
Let $\Lcal$ be a maximal family of cross-free cuts of a graph $G= ([n],w)$.  Then $|\Lcal| = 2n-3$.
\end{corollary}

\begin{proof}
Let $\Gcal$ be the beach of $\Lcal$ and $(T,\phi)$ a faithful tree representation of $\Gcal$.  As $(T,\phi)$
is faithful $|E(T)| = |\Lcal|$.  Let $T'$ be the undirected graph underlying $T$.  Clearly $|E(T')| = |E(T)|$.  We use \cref{lem:tree_structure} to count $|E(T')|$.
Let $i$ be the number of internal vertices of $T'$, each of which has degree 3.  There are also $n$ non-internal vertices each of which has
degree $1$.  Thus $|E(T')| = (3i + n)/2$.  Also as $T'$ is a tree $|E(T')| = |V(T')|-1 = n+i-1$.  Hence $i = n-2$ and $|\Lcal| = |E(T)| = |E(T')| = 2n-3$.
\end{proof}

\subsection{Linear independence}
We now show the main theorem of this section that in a complete weighted graph any set $\vec{\Lcal}$ of cut vectors of a cross-free family of cuts $\Lcal$ is linearly independent.
We will use the tree representation $(T,\phi)$ of the beach $\Gcal$ of $\Lcal$ to do this via the following lemma.
\begin{lemma}
\label{lem:tree_assign}
Let $T$ be an arborescence with root $r$ and $\psi : E(T) \rightarrow \R$.   Let $U$ be a finite set and $\phi: U \rightarrow V(T)$.  Suppose that $T,\phi, \psi$ have the property
that
\begin{enumerate}
  \item The root $r$ is labeled and has $|\delta^+(r)| = 1$.
  \item Every internal vertex $v$ is unlabeled and has $|\delta^+(v)| = 2$.
  \item Every leaf of $T$ has a label.
  \item For every $s,t \in U$ it holds that $\sum_{e \in \phi(s)-\phi(t)} \psi(e) = 0$, where $\phi(s)-\phi(t)$ is the set of edges on the undirected path from $\phi(s)$ to $\phi(t)$.
\end{enumerate}
Then $\psi$ is identically $0$.
\end{lemma}

\begin{proof}
We will prove by induction on the depth of the arborescence.  We need a slightly different statement for the inductive hypothesis
since when considering a sub-arborescence $T'$ of $T$ we do not know that the root of $T'$ has property (1).

\paragraph{Inductive hypothesis:} Let $T$ be an arborescence with root $r$ that is unlabeled and has $|\delta^+(r)| = 2$, and further suppose
$T, \phi, \psi$ satisfy conditions (2)-(4) of the proposition.  Then letting $u,v$ be the children of $r$ it holds that
$\psi((r,u)) = -\psi((r,v))$ and for any other edge $e \in E(T), e \ne (r,u),(r,v)$ it holds that $\psi(e) = 0$.

For the base case consider a tree of depth $1$, with root $r$ and two children $u,v$ which are leaves.  As they are leaves, $u,v$ are
labeled which, considering the path from $u$ to $v$, means $\psi((r,u)) + \psi((r,v)) = 0$.  This concludes the base case.

Now we prove the inductive step.  Let $r$ be the root of a tree with children $u,v$.  We consider two cases:
\paragraph{Case 1: one of $u,v$ is a leaf.} Suppose without loss of generality that $u$ is a leaf and $v$ is an internal node with children $v_1, v_2$.
By the inductive hypothesis $\psi((v,v_1)) + \psi((v,v_2)) = 0$ and $\psi$ is identically $0$ on the subtrees rooted at $v_1,v_2$.  Let $y_1, y_2$ be leaves
that are descendants of $v_1,v_2$ respectively (and can possibly be $y_1,y_2$ themselves).  Considering the path from $u$ to $y_1$ and $y_2$ we have
the equations
\begin{align*}
\psi((r,u) + \psi((r,v)) + \psi((v,v_1)) &= 0 \\
\psi((r,u) + \psi((r,v)) + \psi((v,v_2)) &= 0
\end{align*}
As $\psi((v,v_1)) + \psi((v,v_2)) = 0$, adding these equations shows that $\psi((r,u) + \psi((r,v)) =0$, as desired.  Substituting this back
into the equations further implies that $\psi((v,v_1)) = \psi((v,v_2)) = 0$ so $\psi$ is identically $0$ on the subtree rooted at $v$ completing
this case.

\paragraph{Case 2: both $u,v$ are internal vertices.} Let the children of $u$ be $u_1, u_2$ and the
children of $v$ be $v_1, v_2$.
By the inductive hypothesis, $\psi(\cdot)$ is identically zero on the sub-trees rooted at $u_1, u_2, v_1, v_2$
and we have $\psi((u,u_1)) + \psi((u,u_2)) = \psi((v,v_1)) + \psi((v,v_2))= 0$.  We must show that $\psi((u,u_1))= \psi((u,u_2))=\psi((v,v_1)) =\psi((v,v_2))= 0$
and that $\psi((r,u)) + \psi((r,v)) = 0$.

Let $x_1,x_2$ be a leaves that are descendants of $u_1,u_2$, respectively, and similarly let $y_1, y_2$ be leaves that are
descendants of $v_1,v_2$, respectively.  By assumption all of these leaves are labeled.  Considering the paths from $x_b-y_{b'}$ for $b,b' \in \{0,1\}$ we obtain
the following four constraints on $\psi$:
\begin{align*}
\psi((u,u_1)) + \psi((r,u)) + \psi((r,v)) + \psi((v,v_1)) &= 0 \\
\psi((u,u_1)) + \psi((r,u)) + \psi((r,v)) + \psi((v,v_2)) &= 0 \\
\psi((u,u_2)) + \psi((r,u)) + \psi((r,v)) + \psi((v,v_1)) &= 0 \\
\psi((u,u_2)) + \psi((r,u)) + \psi((r,v)) + \psi((v,v_2)) &= 0 \\
\end{align*}
Adding all four equations and using $\psi((u,u_1)) + \psi((u,u_2)) = \psi((v,v_1)) + \psi((v,v_2))= 0$ shows that $\psi((r,u)) + \psi((r,v)) = 0$.
Taking this into account, adding the first two equations then shows $\psi((u,u_1)) = 0$, and adding the last two equations shows $\psi((u,u_2)) = 0$.
This then also means $\psi((v,v_1)) = \psi((v,v_2)) = 0$.

We have now shown the inductive statement holds.  It remains to see why this implies the lemma.  Let $r$ be the root of the tree, let $u$ be the
child of $r$, and let $u_1, u_2$ be the children of $u$.  By the inductive statement we have that $\psi((u,u_1)) + \psi((u,u_2)) = 0$ and $\psi$ is
identically zero on the subtree rooted at $u_1$ and the subtree rooted at $u_2$.  Let $x_1,x_2$ be leaves which are descendants of $u_1, u_2$, respectively.
As the root has a label, considering the path from $r$ to $u_1$ implies that $\psi((r,u)) + \psi((u,u_1)) = 0$ and considering the path from $r$ to $u_2$
implies $\psi((r,u)) + \psi((u,u_2)) = 0$.  Adding these equations implies that $\psi((r,u)) = 0$, from which it then follows that $\psi((u,u_1)) = \psi((u,u_2)) = 0$.
\end{proof}

\begin{lemma}
\label{prop:lin_ind}
Let $G = ([n],w)$ be a complete weighted graph and let $\Lcal$ be a cross-free family of cuts.  Then $\vec{\Lcal} = \{\chi(S): S \in \Lcal\}$ form a linearly independent set of vectors.
\end{lemma}

\begin{proof}
We may assume that $\Lcal$ is a \emph{maximal} cross-free family, as showing that a superset of $\vec{\Lcal}$ is linearly independent
implies that $\vec{\Lcal}$ is as well.  Thus suppose $\Lcal$ is a maximal cross-free family and let $\Gcal$ be its beach.  Let $(T,\phi)$ be a faithful tree representation of
$\Gcal$.  By \cref{lem:tree_structure} we have that $(T,\phi)$ satisfy conditions (1)-(3) of \cref{lem:tree_assign}.

Now we ask the question: for an edge $\{i,j\} \in E(G)$ which sets $S \in \Lcal$ contain it?
This has a very nice description in terms of the tree decomposition.  Let $u,v \in V(T)$ be the vertices with
$\phi(i) = u, \phi(j) = v$.  Then the sets containing $i$ are the sets represented by edges from the root to $u$;  the sets containing
$j$ are the sets represented by the edges on the path from the root to $v$.  Therefore the sets which contain $i$ but not $j$ or $j$ but
not $i$, are exactly those represented by the edges on the path from $u$ to $v$ in the undirected tree underlying $T$.  Thus the cuts which
contain the edge $\{i,j\}$ are exactly those with a shore which is represented by an edge on the path from $u$ to $v$ in undirected graph underlying $T$.

Consider a linear combination $\sum_{S \in \Lcal} \alpha_S \chi(S) = \zero$ which is equal to the all zero vector.  The $\{i,j\}$ coordinate
of this equation says that $\sum_{S \in \Lcal, \{i,j\} \in S} \alpha_S \chi(S)(\{i,j\}) = 0$.  This sum is exactly over the sets represented by edges
on the path from $\phi(i)$ to $\phi(j)$.  As this sum must be zero for every edge $\{i,j\}$, this says that if we let $\psi(e) = \alpha_S$ where the edge $e$
represents a shore of $S$ then for any two labeled vertices $u,v \in V(T)$ the sum of $\psi(e)$ over the edges on the path from $u$ to $v$ is zero.
Thus also condition~(4) of \cref{lem:tree_assign} is satisfied.
Hence all of the conditions of \cref{lem:tree_assign} hold which implies that
$\psi$ must be identically zero and therefore all coefficients $\alpha_S = 0$.  This shows that $\{\chi(S): S \in \Lcal\}$ is a linearly independent set.
\end{proof}

We can now give the first proof of our main lower bound result on the cut dimension \cref{thm:cd_lower}, which
says that for every integer $n\ge 2$ there is an $n$-vertex weighted graph $G = (V,w)$ with $\cd(G) \ge 2n-3$.

\begin{proof}[Proof of \cref{thm:cd_lower}]
For every integer $n \ge 2$, Theorem~5.2 of \cite{CR04} constructs a complete weighted graph $G = (V,w)$ on $n$ vertices such that
$\M(G)$ is a cross-free family of size $|\M(G)| = 2n-3$.  By \cref{prop:lin_ind} the vectors in $\vec{\Mcal}$ form a linearly
independent set, thus $\cd(G) \ge 2n-3$.
\end{proof}

\subsection{Constructing graphs with a cross-free set of mincuts}
\label{subsec:construction}
In this subsection we explicitly construct, for any maximal cross-free family $\Fcal \subseteq 2^{[n]}$, a complete weighted graph $G = ([n],w)$
with $\M(G) = \{ \Delta(X) : X \in \Fcal\}$.  This task is made easier by the next lemma.  We first need a definition.

\begin{definition}
\label{def:overlap}
Let $\Fcal \subseteq 2^{V}$.  For a subset $X \subseteq V$, let $\cross_{\Fcal}(X) = \{Y \in \Fcal: X, Y \ \cross\}$.
\end{definition}

\begin{lemma}
\label{lem:min_value}
Let $G = (V,w)$ be a graph and $\Lcal$ be a maximal cross-free family of cuts.  Suppose that for all
$S \in \Lcal$ it holds that $w(S) = c$.  Then the weight of a minimum cut in $G$ is $c$.
\end{lemma}

\begin{proof}
Let $\Gcal$ be the beach of $\Lcal$.  Suppose for a contradiction that the weight of a minimum cut of $G$ is $< c$.
Let $\Tcal = \{ Z : \emptyset \ne Z \subsetneq V, v_1 \not \in Z, Z \not \in \Gcal, w(\Delta(Z)) < c\}$ and
\[
X = \argmin_{Z} \{|\cross_\Gcal(Z)|: Z \in \Tcal \} \enspace .
\]
In the following we always use $\cross(\cdot)$ with respect to $\Gcal$ and
drop the subscript.  As $|\cross(X)| \ge 1$, let $Y \in \cross(X)$.  As shown in \cref{sec:appendix} \cref{lem:intersect}, both
$|\cross(X \cap Y)|$ and $|\cross(X \cup Y)|$ are strictly smaller than $|\cross(X)|$.  Thus it must be the case that
$X \cap Y, X \cup Y \not \in \Tcal$.  Let us take the case of $X \cap Y$.  It does not contain $v_1$, as neither $X$ nor $Y$ do,
and it is a nonempty set by the definition of $\cross$.  Thus it must be the case that either $w(\Delta(X \cap Y)) \ge c$ or
that $X \cap Y \in \Gcal$, which implies $w(\Delta(X \cap Y)) = c$.  The same argument holds for $X \cup Y$, thus
both $w(\Delta(X \cap Y)), w(\Delta(X \cup Y)) \ge c$.

However by submodularity of the cut function we have
$w(\Delta(X\cap Y )) + w(\Delta(X \cup Y)) \le w(\Delta(X)) + w(\Delta(Y))$, which implies that at least one of $\Delta(X \cap Y), \Delta(X \cup Y)$ must have weight $<c$.
Hence we have a contradiction and the lemma holds.
\end{proof}

We will additionally need the following theorem which follows from Theorem 5.1 in \cite{CR04}.
\begin{theorem}[\cite{CR04}]
\label{thm:cr_complete}
Let $G = (V,w)$ be a complete weighted graph.  Then $\M(G)$ is a cross-free family of cuts.
\end{theorem}


\begin{theorem}
\label{thm:explicit}
Let $n \ge 2$ and $\Lcal$ be a {maximal} cross-free family of cuts in the $n$-vertex complete graph.  Let
$A$ be an $|\Lcal|$-by-$\binom{n}{2}$ matrix whose rows are the vectors $\chi(S)$ for $S \in \Lcal$ and
let $z = A^T \1$.  Define $w(e) = 2^{-z(e)+1}$ for $e \in [n]^{(2)}$.
Then $G=([n],w)$ is a complete weighted graph with $\cd(G) = 2n-3$ and $\M(G) = \Lcal$.
\end{theorem}

\begin{proof}
It is clear from the definition that $w >0$ and so defines a complete weighted graph.
We will show that $Aw = \1$.  By \cref{lem:min_value} this shows that
the minimum cut weight of $G$ is $1$ and so the set of minimum cuts includes $\Lcal$.  As $w$ defines a complete weighted graph, by \cref{thm:cr_complete}
the set of minimum cuts in $G$ is cross-free and therefore must be exactly $\M(G) = \Lcal$, since $\Lcal$ is maximal.  Further, {$|\Lcal| = 2n-3$ by \cref{cor:maxcard} and}
the vectors in $\vec{\Lcal}$ are linearly independent by \cref{prop:lin_ind}, thus $\cd(G) = 2n-3$.

It remains to show $Aw = \1$.
We do this using an alternative way of viewing the assignment of edge weights.
Let $\Gcal \subseteq 2^{[n]}$ be the beach of $\Lcal$, and $(T,\phi)$ be a faithful
tree representation of $\Gcal$.  For vertices $u,v \in V(T)$ let $d(u,v)$ be the length of the shortest path between $u,v$ in the undirected graph underlying $T$.
Now let $\{i,j\} \in [n]^{(2)}$ and suppose $\phi(i) = u, \phi(j) = v$.  We claim that $w(\{i,j\}) = 2^{-d(u,v)+1}$.
The sets of $\Gcal$ containing $i$ are the sets represented by edges from the root to $u$;  the sets of $\Gcal$ containing
$j$ are the sets represented by the edges on the path from the root to $v$.  Therefore the sets which contain $i$ but not $j$ or $j$ but
not $i$, are exactly those represented by the edges on the path from $u$ to $v$ in the undirected tree underlying $T$.  As $(T,\phi)$ is faithful,
each of these edges represents a different set, and therefore the number of edges on the path from $u$ to $v$ is exactly the number of
sets of $\Lcal$ which contain $\{i,j\}$.

We now continue with the proof that $Aw = \1$ using this interpretation of the weights.  For any cut $S\in \Lcal$ with shore $X \in \Gcal$, take the edge $(u,v) \in E(T)$ representing $X$.
Now imagine we remove the edge $(u,v)$ from $T$ which disconnects $T$ into two components.  Let $T_u$ be the component containing $u$ and $T_v$ the component containing $v$.
From $T_u$, which contains the root $r$ of $T$, we create a graph $T_u'$ whose underlying undirected graph is the same as $T_u$, but for which all edges are directed away from $u$.  Thus in
$T_u'$, vertex $u$ becomes the root and $r$ becomes a leaf.  Now by item~(2) of \cref{lem:tree_structure}, every non-leaf vertex in $T_v$ and $T_u'$ has out-degree $2$.  We inject a unit of flow into $u$ in the graph $T_u'$
and let it propagate according to the rule that at every non-leaf vertex half of the flow is routed along each outgoing edge.  We similarly inject a unit of flow into $v$ in the graph $T_v$ and let it
propagate according to the same rule.  Thus in the tree $T_v$, each leaf $a$ gets $f(a) = 2^{-d(a,v)}$ amount of flow, where $d(a,v)$ is the number of edges along the path from $v$ to $a$ in $T_v$.
Similarly, if $b$ is a leaf in the tree $T_u'$, the amount of flow arriving at $b$ is $f(b) = 2^{-d(b,u)}$. Now let $\{i,j\} \in [n]^{(2)}$ with $i \in X, j \in \bar{X}$ and observe that the way we defined $w(\{i,j\})$ satisfies
 \[
 w(\{i,j\}) = 2^{-d(\phi(i),\phi(j))+1}  = 2^{-d(\phi(i),v) - d(\phi(j),u)} = f(\phi(i))\cdot  f(\phi(j)) \enspace.
 \]
 Thus the weight of the cut $S$ is
 \[
 \sum_{i\in X,j\in \bar{X}} w(\{i,j\}) =  \sum_{i\in X,j\in \bar{X}} f(\phi(i)) \cdot f(\phi(j)) =  \left(\sum_{i\in X} f(\phi(i))\right) \cdot \left(\sum_{j\in \bar{X}} f(\phi(j))\right) = 1\cdot 1 = 1 \enspace .
 \]

\end{proof}

%


\appendix\setcounter{theorem}{44}\section{Jain's spanning lemma}
\label{sec:appendix}
In this appendix we prove \cref{thm:span}.  The proof uses the following key property of mincuts which goes back at least
to work of Dinitz, Karzanov, and Lomonosov \cite{DKL76}.

\begin{proposition}[\cite{DKL76} ``Lemma on a quadrangle'']
\label{prop:bixby}
Let $G = (V,w)$ be a graph.
For any crossing mincuts $\Delta(X), \Delta(Y)$ of $G$
it holds that
\[
\chi(\Delta(X)) + \chi(\Delta(Y)) = \chi(\Delta(X \cap Y)) + \chi(\Delta(X \cup Y)) \enspace .
\]
\end{proposition}

\begin{proof}
If $\Delta(X), \Delta(Y)$ cross then
$\Delta(X \cap Y), \Delta(X \cup Y)$ are mincuts of $G$ by \cref{claim:mincuts}.
Further, by counting the number of times an edge appears on each side it can be seen (eg.\ Ex. 6.48 in \cite{Lov93}) that
\begin{equation}
\label{eq:cut_vecs}
\chi(\Delta(X)) + \chi(\Delta(Y)) = \chi(\Delta(X \cap Y)) + \chi(\Delta(X \cup Y)) + 2\chi(E(X-Y, Y-X)) \enspace.
\end{equation}

Let the minimum cut value of $G$ be $\lambda$.  Let $m$ be the number of edges in $G$ and
$\vec{w} \in \R^m$ be the positive vector resulting from restricting $w$ to the edges of $G$.
The inner product of $\vec{w}$ with the left hand side of \cref{eq:cut_vecs}
is $2\lambda$, and with the righthand side is $2\lambda  + 2 \braket{\vec{w}}{\chi(E(X-Y, Y-X))}$.  Thus $\braket{\vec{w}}{\chi(E(X-Y, Y-X))} = 0$,
which implies $\chi(E(X-Y, Y-X)) = \zero$ since
$\vec{w}$ is positive and $\chi(E(X-Y, Y-X))$ is nonnegative.
\end{proof}


Jain's proof uses the technique of \emph{combinatorial uncrossing}.
Recall the definition of $\cross_{\Gcal}(X)$ from \cref{def:overlap}.   A key to the proof is the following simple lemma about
$\cross_{\Gcal}(X)$.
%
\begin{lemma}[\cite{Jain01}]
\label{lem:intersect}
Let $\Fcal \subseteq 2^{V}$ be a set family closed under overlaps and $\Gcal \subseteq \Fcal$ be a maximal laminar subset of $\Fcal$.
Then for any $X \in \Fcal - \Gcal$ and $Y \in \cross_{\Gcal}(X)$
\begin{align}
\label{eq:con_cap}
\cross_{\Gcal}(X \cap Y) & \subset \cross_{\Gcal}(X) \\
\label{eq:con_cup}
\cross_{\Gcal}(X \cup Y) &\subset \cross_{\Gcal}(X) \enspace .
\end{align}
\end{lemma}


\begin{proof}
In the following we always refer to $\cross(X)$ with respect to $\Gcal$ and drop the subscript.
We first show \cref{eq:con_cap}.  First note that $Y \in \cross(X) - \cross(X \cap Y)$.  Thus to show \cref{eq:con_cap} it suffices to show $\cross(X \cap Y) \subseteq \cross(X)$.
Let $W \in \cross(X \cap Y)$.  We want to show that $W \in \cross(X)$, i.e.\ that it cannot be the case that $W \subseteq X, X \subseteq W$, or
$X \cap W = \emptyset$.  We know that the last one cannot hold as $W \cap (X \cap Y) \ne \emptyset$ as $W \in \cross(X \cap Y)$.

Also as $W,Y \in \Gcal$ they do not overlap and thus either $Y \subseteq W, W \subseteq Y$, or $Y \cap W = \emptyset$.  Again the last one
cannot hold as $W \cap (X \cap Y) \ne \emptyset$.  The following table shows that assuming $W \not \in \cross(X)$ leads to a contradiction
in all $4$ remaining cases.
\begin{center}
\begin{tabular}{c|c|c|}
                                       &      $Y \subseteq W$       & $W \subseteq Y$                              \\
\hline
$W \subseteq X$
& \begin{tabular}{@{}c@{}} $Y \subseteq X$ \\ $Y \not \in \cross(X)$ \end{tabular}
& \begin{tabular}{@{}c@{}} $W \subseteq X \cap Y$ \\ $W \not \in \cross(X \cap Y)$ \end{tabular}   \\
\hline
$X \subseteq W$
&  \begin{tabular}{@{}c@{}} $X \cap Y \subseteq W$ \\ $W \not \in \cross(X \cap Y)$ \end{tabular}
& \begin{tabular}{@{}c@{}} $X \subseteq Y$ \\ $Y \not \in \cross(X)$ \end{tabular}  \\
\hline
\end{tabular}
\end{center}

We now show \cref{eq:con_cup}, which follows similarly.  Again $Y \in \cross(X) - \cross(X \cup Y)$ thus it suffices to show
$\cross(X \cup Y) \subseteq \cross(X)$.  Let $W \in \cross(X \cup Y)$.  We want to show that $W \in \cross(X)$, i.e.\ that
is not the case that either $W \cap X = \emptyset, X \subseteq W$, or $W \subseteq X$.  We cannot have $W \subseteq X$
because this means $W \subseteq X \cup Y$ which contradicts $W \in \cross(X \cup Y)$.  As $W,Y \in \Gcal$ they do
not overlap, so we also know either $Y \subseteq W, W \cap Y = \emptyset$, or $W \subseteq Y$.  The last one again
cannot hold as it implies $W \subseteq X \cup Y$.  The following table shows that assuming $W \not \in \cross(X)$ leads
to a contradiction in the remaining $4$ cases.
\begin{center}
\begin{tabular}{c|c|c|}
                                       &      $Y \subseteq W$            & $W \cap Y = \emptyset$ \\
\hline
$X \subseteq W$
& \begin{tabular}{@{}c@{}} $X \cup Y \subseteq W$ \\ $W \not \in \cross(X \cup Y)$ \end{tabular}
& \begin{tabular}{@{}c@{}} $X \cap Y = \emptyset$ \\ $Y \not \in \cross(X)$ \end{tabular} \\
\hline
$X \cap W = \emptyset$
&  \begin{tabular}{@{}c@{}} $Y \cap X = \emptyset$ \\ $Y \not \in \cross(X)$ \end{tabular}
& \begin{tabular}{@{}c@{}} $W \cap (X \cup Y) = \emptyset$ \\ $W \not \in \cross(X \cup Y)$ \end{tabular} \\
\hline
\end{tabular}
\end{center}
\end{proof}

We are now ready to show the key lemma of Jain.
\jainspan*

\begin{proof}
It is clear that $\laspan(\vec{\Lcal}) \subseteq \laspan(\vec{\M}(G))$ so we focus on the other direction.

Let $\Fcal$ be the beach of $\M(G)$.  By \cref{claim:mincuts} $\Fcal$ is closed under overlaps.  Let $\Gcal \subseteq \Fcal$ be the beach of $\Lcal$.  As $\Lcal$ is a
maximal cross-free subset of $\M(G)$ it follows that $\Gcal$ is a maximal laminar subset of $\Fcal$.
Thus $|\cross_{\Gcal}(X)| \ge 1$ for all $X \in \Fcal - \Gcal$.  In the following we will always refer to $\cross(X)$ with respect to $\Gcal$ and drop the subscript.

Suppose for a contradiction that $\laspan(\vec{\Lcal}) $ is a strict subset of
$\laspan(\vec{\Mcal}(G))$.  Let
\[
X = \argmin_{Z \in \Fcal - \Gcal} \{ |\cross(Z) |:  \chi(\Delta(Z)) \not \in  \laspan(\vec{\Lcal}) \} \enspace .
\]
As $\cross(X) \ge 1$, let $Y \in \cross(X)$.  By \cref{lem:intersect}
\begin{align}
\label{eq:card_cap}
|\cross(X \cap Y)| &< |\cross(X)| \\
\label{eq:card_cup}
|\cross(X \cup Y)| &< |\cross(X)| \enspace .
\end{align}
By the definition of $X$, and as $\Fcal$ is closed under overlaps, we must have
$\chi(\Delta(X \cap Y)), \chi(\Delta(X \cup Y)) \in \laspan(\vec{\Lcal})$.
Also as $Y \in \Gcal$ we have $\chi(\Delta(Y)) \in \vec{\Lcal}$ which implies by \cref{prop:bixby} that
\[
\chi(\Delta(X)) =  \chi(\Delta(X \cap Y))+ \chi(\Delta(X \cup Y)) - \chi(\Delta(Y)) \enspace.
\]
This implies $\chi(\Delta(X)) \in \laspan(\vec{\Lcal})$, a contradiction.
\end{proof}

\begin{thebibliography}{GMW20b}

\bibitem[ACK20]{ACK20}
Sepehr Assadi, Deeparnab Chakrabarty, and Sanjeev Khanna.
\newblock Graph connectivity and single element recovery via linear and {OR}
  queries.
\newblock {\em CoRR}, abs/2007.06098, 2020.

\bibitem[BFS86]{BFS86}
L{\'{a}}szl{\'{o}} Babai, Peter Frankl, and Janos Simon.
\newblock Complexity classes in communication complexity theory (preliminary
  version).
\newblock In {\em 27th Annual Symposium on Foundations of Computer Science,
  Toronto, Canada, 27-29 October 1986}, pages 337--347, 1986.

\bibitem[BG08]{BG08}
Andr{\'{a}}s~A. Bencz{\'{u}}r and Michel~X. Goemans.
\newblock Deformable polygon representations and near-mincuts.
\newblock In Martin Gr{\o}tschel and Gyula O.~H. Katona, editors, {\em Building
  Bridges: Between Mathematics and Computer Science}, volume~19 of {\em Bolyai
  Society Mathematical Studies}, pages 103--135. Springer, 2008.

\bibitem[Bix75]{bixby75}
R.~E. Bixby.
\newblock The minimum number of edges and vertices in a graph with edge
  connectivity n and m n-bonds.
\newblock {\em Netw.}, 5(3):253–298, July 1975.

\bibitem[CR04]{CR04}
L.~Sunil Chandran and L.~Shankar Ram.
\newblock On the number of minimum cuts in a graph.
\newblock {\em {SIAM} J. Discret. Math.}, 18(1):177--194, 2004.

\bibitem[DKL76]{DKL76}
Efim~A. Dinitz, Alexander~V. Karzanov, and Michael~V. Lomonosov.
\newblock On the structure of the system of minimum edge cuts of a graph.
\newblock {\em Studies in discrete optimization}, 1976.

\bibitem[FF09]{FF09}
Tam{\'{a}}s Fleiner and Andr{\'{a}}s Frank.
\newblock A quick proof for the cactus representation of mincuts.
\newblock {\em EGRES Quick Proof}, 2009-03, 2009.

\bibitem[GMW20a]{GMW20}
Pawel Gawrychowski, Shay Mozes, and Oren Weimann.
\newblock Minimum cut in ${O}(m \log^2 n)$ time.
\newblock In Artur Czumaj, Anuj Dawar, and Emanuela Merelli, editors, {\em 47th
  International Colloquium on Automata, Languages, and Programming, {ICALP}
  2020, July 8-11, 2020, Saarbr{\"{u}}cken, Germany (Virtual Conference)},
  volume 168 of {\em LIPIcs}, pages 57:1--57:15. Schloss Dagstuhl -
  Leibniz-Zentrum f{\"{u}}r Informatik, 2020.

\bibitem[GMW20b]{GMW20b}
Pawel Gawrychowski, Shay Mozes, and Oren Weimann.
\newblock A note on a recent algorithm for minimum cut.
\newblock {\em CoRR}, abs/2008.02060, 2020.

\bibitem[GNT20]{GNT20}
Mohsen Ghaffari, Krzysztof Nowicki, and Mikkel Thorup.
\newblock Faster algorithms for edge connectivity via random 2-out
  contractions.
\newblock In Shuchi Chawla, editor, {\em Proceedings of the 2020 {ACM-SIAM}
  Symposium on Discrete Algorithms, {SODA} 2020, Salt Lake City, UT, USA,
  January 5-8, 2020}, pages 1260--1279. {SIAM}, 2020.

\bibitem[Goe06]{Goemans06}
Michel~X. Goemans.
\newblock Minimum bounded degree spanning trees.
\newblock In {\em 47th Annual {IEEE} Symposium on Foundations of Computer
  Science {(FOCS} 2006), 21-24 October 2006, Berkeley, California, USA,
  Proceedings}, pages 273--282. {IEEE} Computer Society, 2006.

\bibitem[GPRW20]{GPRW20}
Andrei Graur, Tristan Pollner, Vidhya Ramaswamy, and S.~Matthew Weinberg.
\newblock New query lower bounds for submodular function minimization.
\newblock In Thomas Vidick, editor, {\em 11th Innovations in Theoretical
  Computer Science Conference, {ITCS} 2020, January 12-14, 2020, Seattle,
  Washington, {USA}}, volume 151 of {\em LIPIcs}, pages 64:1--64:16. Schloss
  Dagstuhl - Leibniz-Zentrum f{\"{u}}r Informatik, 2020.

\bibitem[GR95]{GR95}
Michel~X. Goemans and V.~S. Ramakrishnan.
\newblock Minimizing submodular functions over families of sets.
\newblock {\em Comb.}, 15(4):499--513, 1995.

\bibitem[Har08]{Harvey08}
Nicholas J.~A. Harvey.
\newblock {\em Matchings, matroids and submodular functions}.
\newblock PhD thesis, Massachusetts Institute of Technology, Cambridge, MA,
  {USA}, 2008.

\bibitem[HMT88]{HMT88}
Andr{\'{a}}s Hajnal, Wolfgang Maass, and Gy{\"{o}}rgy Tur{\'{a}}n.
\newblock On the communication complexity of graph properties.
\newblock In {\em Proceedings of the 20th Annual {ACM} Symposium on Theory of
  Computing, May 2-4, 1988, Chicago, Illinois, {USA}}, pages 186--191, 1988.

\bibitem[HRW20]{HRW20}
Monika Henzinger, Satish Rao, and Di~Wang.
\newblock Local flow partitioning for faster edge connectivity.
\newblock {\em {SIAM} J. Comput.}, 49(1):1--36, 2020.

\bibitem[HW96]{HW96}
Monika~Rauch Henzinger and David~P. Williamson.
\newblock On the number of small cuts in a graph.
\newblock {\em Inf. Process. Lett.}, 59(1):41--44, 1996.

\bibitem[Jai01]{Jain01}
Kamal Jain.
\newblock A factor 2 approximation algorithm for the generalized {S}teiner
  network problem.
\newblock {\em Comb.}, 21(1):39--60, 2001.

\bibitem[Kar93]{Karger93}
David~R. Karger.
\newblock Global min-cuts in {RNC}, and other ramifications of a simple min-cut
  algorithm.
\newblock In Vijaya Ramachandran, editor, {\em Proceedings of the Fourth Annual
  {ACM/SIGACT-SIAM} Symposium on Discrete Algorithms, 25-27 January 1993,
  Austin, Texas, {USA}}, pages 21--30. {ACM/SIAM}, 1993.

\bibitem[Kar00]{Karger00}
David~R. Karger.
\newblock Minimum cuts in near-linear time.
\newblock {\em J. {ACM}}, 47(1):46--76, 2000.

\bibitem[KT19]{KT19}
Ken{-}ichi Kawarabayashi and Mikkel Thorup.
\newblock Deterministic edge connectivity in near-linear time.
\newblock {\em J. {ACM}}, 66(1):4:1--4:50, 2019.

\bibitem[KV18]{KV18}
Bernhard Korte and Jens Vygen.
\newblock {\em Combinatorial Optimization: Theory and Algorithms}.
\newblock Springer, 2018.

\bibitem[Lov93]{Lov93}
L{\'{a}}szl{\'{o}} Lov{\'{a}}sz.
\newblock {\em Combinatorial problems and exercises {(2.} ed.)}.
\newblock North-Holland, 1993.

\bibitem[MN20]{MN20}
Sagnik Mukhopadhyay and Danupon Nanongkai.
\newblock Weighted min-cut: sequential, cut-query, and streaming algorithms.
\newblock In {\em Proceedings of the 52nd Annual ACM SIGACT Symposium on Theory
  of Computing, STOC 2020, Chicago, IL, USA, June 22-26, 2020}, pages 496--509,
  2020.

\bibitem[NNI97]{NNI97}
Hiroshi Nagamochi, Kazuhiro Nishimura, and Toshihide Ibaraki.
\newblock Computing all small cuts in an undirected network.
\newblock {\em {SIAM} J. Discret. Math.}, 10(3):469--481, 1997.

\bibitem[RSW18]{RSW18}
Aviad Rubinstein, Tselil Schramm, and S.~Matthew Weinberg.
\newblock Computing exact minimum cuts without knowing the graph.
\newblock In {\em 9th Innovations in Theoretical Computer Science Conference,
  {ITCS} 2018, January 11-14, 2018, Cambridge, MA, {USA}}, pages 39:1--39:16,
  2018.

\end{thebibliography}
\end{document}
